\documentclass{article}
\usepackage{amsmath, amssymb, amsthm}
\usepackage{hyperref}
\usepackage{geometry}
\geometry{margin=1in}

\title{Erd\H{o}s Problem \#919}
\date{Last updated: 2025-08-31}
\author{Source: erdosproblems.com}

\begin{document}
\maketitle

\noindent\textbf{Status:} OPEN \quad \textbf{No prize}

\noindent\textbf{Tags:} graph theory, chromatic number

\medskip

\noindent\textbf{Problem Statement:}

Is there a graph $G$ with vertex set $\omega_2^2$ and chromatic number $\aleph_2$ such that every subgraph whose vertices have a lesser type has chromatic number $\leq \aleph_0$? What if instead we ask for $G$ to have chromatic number $\aleph_1$?




This question was inspired by a theorem of Babai, that if $G$ is a graph on a well-ordered set with chromatic number $\geq \aleph_0$ there is a subgraph on vertices with order-type $\omega$ with chromatic number $\aleph_0$. Erd\H{o}s and Hajnal showed this does not generalise to higher cardinals - they (see \cite{Er69b}) constructed a set on $\omega_1^2$ with chromatic number $\aleph_1$ such that every strictly smaller subgraph has chromatic number $\leq \aleph_0$ as follows: the vertices of $G$ are the pairs $(x_\alpha,y_\beta)$ for $1\leq \alpha,\beta &#60;\omega_1$, ordered lexicographically. We connect $(x_{\alpha_1},y_{\beta_1})$ and $(x_{\alpha_2},y_{\beta_2})$ if and only if $\alpha_1&#60;\alpha_2$ and $\beta_1&#60;\beta_2$. A similar construction produces a graph on $\omega_2^2$ with chromatic number $\aleph_2$ such that every smaller subgraph has chromatic number $\leq \aleph_1$.




References

[Er69b] Erd\H{o}s, P., Problems and results in chromatic graph theory . Proof Techniques in Graph Theory (Proc. Second Ann Arbor Graph Theory Conf., Ann Arbor, Mich., 1968) (1969), 27-35.


\bigskip
\noindent\small{Source: \url{https://www.erdosproblems.com/919}}
\end{document}
